% संपूर्ण मराठी ग्रंथातील प्रकरण 31: अपूर्णता आणि सिद्धतायोग्यता
% हा संपादनयोग्य स्रोतखंड आहे. संकलनासाठी संपूर्ण स्रोतसंग्रहातील
% अक्षररूपे, पूर्वभाग, चिन्हव्याख्या आणि आकृती-स्रोत आवश्यक आहेत.
% मूळ: Open Logic Project; मजकूर आणि रूपांतर CC BY 4.0.
% यंत्रानुवाद, दुरुस्ती आणि तपासणी: OpenAI Codex —
% GPT-5.6 Sol आणि GPT-6 Sol, दोन्ही Ultra effort.
% दुरुस्ती-पुनरावलोकन: GPT-6.1 Sol, Ultra effort.
\chapter{अपूर्णता आणि सिद्धतायोग्यता}\label{inc:inp:chap}









\section{परिचय}\label{inc:inp:int:sec}

हिल्बर्टच्या मते, नैसर्गिक संख्यांसारख्या गणितीय रचनेसाठीची
स्वयंसिद्धकपद्धती त्या रचनेविषयीची सर्व सत्य विधाने सिद्ध करू देत
नसेल, तर ती अपुरी आहे. निष्पत्तीच्या आकारिक पद्धतींविषयीच्या
त्याच्या नंतरच्या स्वारस्याचा विचार केला, तर नैसर्गिक संख्यांविषयी
युक्तिवाद करण्यासाठी वापरली जाणारी आकारिक पद्धत केवळ सुसंगतच
नव्हे, तर \emph{संपूर्ण}ही असावी, असे त्याला वाटत होते असे दिसते:
त्या भाषेतील प्रत्येक विधान किंवा त्याचा निषेध निष्पन्न करता येण्याजोगे
असावा. गोडेलचे पहिले अपूर्णता प्रमेय दाखवते की अशी
स्वयंसिद्धकपद्धती अस्तित्वात नाही: अंकगणितासाठी संपूर्ण, सुसंगत
आणि स्वयंसिद्धकीकरणीय अशी कोणतीही आकारिक पद्धत नाही. खरे तर,
“पुरेशी प्रबळ”, सुसंगत आणि स्वयंसिद्धकीकरणीय अशी कोणतीही
गणितीय उपपत्ती संपूर्ण नसते.

हिल्बर्टचा आणखी महत्त्वाचा उद्देश म्हणजे अभिजात गणिताला आधार
देण्याच्या कार्यक्रमाचा मुख्य भाग: अभिजात युक्तिवादाचे निरूपण
करणाऱ्या आकारिक पद्धतींच्या सुसंगततेच्या सांतवादी सिद्धता
शोधणे. म्हणून हिल्बर्टच्या कार्यक्रमासाठी गोडेलचे दुसरे
अपूर्णता प्रमेय हा अधिक मोठा धक्का होता. पहिल्या प्रमेयाप्रमाणे
दुसरेही प्रमेय काहीसे स्थूल शब्दांत सांगता येते. साधारणपणे ते
असे म्हणते की अंकगणिताची कोणतीही सुसंगत, प्रभावीपणे
स्वयंसिद्धकीकरण करता येणारी आणि पुरेशी प्रबळ उपपत्ती स्वतःची
सुसंगतता व्यक्त करणारे प्रमाण आकारिक वाक्य सिद्ध करू शकत नाही. येथे “पुरेशी प्रबळ” यामध्ये
केवळ~$\Th{Q}$ पेक्षा थोडे अधिक सामर्थ्य समाविष्ट करावे लागेल.

गोडेलच्या अपूर्णता प्रमेयाच्या मूळ सिद्धतेमागची कल्पना
एपिमेनिडीजच्या विरोधाभासात दिसते. क्रीटचा रहिवासी असलेल्या
एपिमेनिडीजने सर्व क्रीटवासी खोटारडे असतात असे म्हटले;
“हे वाक्य असत्य आहे” हे विधान त्या विरोधाभासाचे अधिक थेट
रूप आहे. सारतः, सत्याच्या जागी निष्पन्नता ठेवून गोडेलने
स्वतःच निष्पन्न करता येण्याजोगे नसल्याचे अप्रत्यक्षपणे सांगणारे
वाक्य आकारिक रीतीने तयार केले. ते वाक्य
निष्पन्न करता येण्याजोगे असेल, तर उपपत्ती विसंगत ठरेल. गोडेलने
विचारात घेतलेल्या स्वयंसिद्धकपद्धतीत त्या वाक्याचा
निषेधही निष्पन्न करता येण्याजोगे नसतो, असे त्याने दाखवले. (या दुसऱ्या
भागासाठी त्याला उपपत्ती~$\Th{T}$ ही तथाकथित
“$\omega$-सुसंगत” आहे असे गृहीत धरावे लागले.
$\omega$-सुसंगतता ही सुसंगततेशी संबंधित पण अधिक प्रबळ अट
आहे.\footnote{म्हणजे प्रत्येक $\omega$-सुसंगत उपपत्ती
सुसंगत असते; उलट विधान खरे नसते.} गोडेलनंतर काही वर्षांनी
रोसरने केवळ~$\Th{T}$ च्या साध्या सुसंगततेचे गृहीतक पुरेसे
आहे असे दाखवले.)

पहिली अडचण म्हणजे स्वतःचा निर्देश करणारे वाक्य
कसे तयार करायचे हे समजणे. $\Th{Q}$ च्या भाषेतील प्रत्येक
सूत्र~$\varphi$ साठी~$\gn{\varphi}$ हा~$\Gn{\varphi}$ या संख्येला
अनुरूप संख्यांक मानू. याचा अर्थ नीट पाहू: $\varphi$ हे
$\Th{Q}$ च्या भाषेतील सूत्र आहे, $\Gn{\varphi}$ ही
नैसर्गिक संख्या आहे, आणि $\gn{\varphi}$ हे $\Th{Q}$ च्या
भाषेतील \emph{पद} आहे. त्यामुळे $\Th{Q}$ च्या भाषेतील
प्रत्येक सूत्र~$\varphi$ ला~$\gn{\varphi}$ असे \emph{नाव}
मिळते; ते $\Th{Q}$ च्या भाषेतील पद आहे. यामुळे
$\Th{Q}$ च्या भाषेतील सूत्रे ना
इतर सूत्रे विषयी काहीतरी “सांगता” येते अशी
संकल्पनात्मक चौकट मिळते. पुढील उपपत्तीला स्थिरबिंदू
उपपत्ती म्हणतात.

\begin{lem}
$\Th{T}$ ही~$\Th{Q}$ चा विस्तार असलेली कोणतीही उपपत्ती
असू द्या, आणि $\psi(x)$ हे फक्त~$x$ हेच मुक्त चल असलेले
कोणतेही सूत्र असू द्या. मग असे
वाक्य~$\varphi$ असते की
$\Th{T} \Proves\varphi \liff \psi(\gn{\varphi})$.
\end{lem}

ही उपपत्ती असे सांगते की कोणत्याही गुणधर्मासाठी $\psi(x)$,
“$\psi(x)$ माझ्याबाबतीत खरा आहे” असे सांगणारे
वाक्य~$\varphi$ असते, आणि $\Th{T}$ ला हे
“माहीत” असते.

असे वाक्य कसे तयार करायचे? क्वाइनने मांडलेले
एपिमेनिडीजच्या विरोधाभासाचे पुढील रूप पाहू:
\begin{quote}
“उद्धृत रूप आधी जोडल्यास असत्य विधान मिळते” उद्धृत रूप
आधी जोडल्यास असत्य विधान मिळते.
\end{quote}
हे वाक्य थेट स्वतःचा निर्देश करत नाही. ते फक्त अवतरणचिन्हांतील
विन्यासात्मक वस्तूंविषयी विधान करते; त्या अर्थाने ते
पुढील वाक्यांसारखेच आहे:
\begin{enumerate}
\item “रॉबर्ट” हे छान नाव आहे.
\item “मी धावलो.” हे छोटे वाक्य आहे.
\item “तीन शब्द आहेत” यात तीन शब्द आहेत.
\end{enumerate}
पण “उद्धृत रूप आधी जोडल्यास असत्य विधान मिळते” या
वाक्यांशाच्या आधी त्याचेच उद्धृत रूप जोडले, तर काय होते?
मूळ वाक्यच मिळते!{} थोडक्यात, ते वाक्य स्वतःच असत्य
असल्याचे सांगते.












\section{स्थिरबिंदू उपपत्ती}\label{inc:inp:fix:sec}

\begin{explain}
स्थिरबिंदू उपपत्ती सांगते की कोणत्याही
सूत्र~$\psi(x)$ साठी असे वाक्य~$\varphi$ असते की
$\Th{T} \Proves \varphi \liff \psi(\gn{\varphi})$, मात्र
$\Th{T}$ हा~$\Th{Q}$ चा विस्तार असला पाहिजे.
खोटारड्याच्या वाक्यासाठी $\varphi$ हे~“$\gn{\varphi}$ असत्य आहे”
या विधानाशी (उपपत्ती~$\Th{T}$ मध्ये सिद्ध होईल अशा
रीतीने) सममूल्य हवे; म्हणजेच $\Gn{\varphi}$ ही एका असत्य
वाक्याची गोडेल संख्या आहे असे त्यात म्हटलेले हवे.
सिद्धतेची कल्पना समजून घेण्यासाठी $\varphi$ विषयी क्वाइनच्या
अनौपचारिक मांडणीशी तुलना करू: “उद्धृत रूप आधी जोडल्यास
असत्य विधान मिळते” उद्धृत रूप आधी जोडल्यास असत्य विधान
मिळते. एखादा शब्दसमूह घेऊन त्याच्या आधी त्याचेच उद्धृत
रूप जोडून वाक्य तयार करण्याच्या क्रियेला त्या शब्दसमूहाचे
\emph{विकर्णीकरण} म्हणू शकतो; तयार झालेला परिणाम हे
त्याचे विकर्णरूप. म्हणून “उद्धृत रूप आधी जोडल्यास असत्य
विधान मिळते” या शब्दसमूहाचे विकर्णरूप म्हणजे “उद्धृत रूप
आधी जोडल्यास असत्य विधान मिळते” उद्धृत रूप आधी
जोडल्यास असत्य विधान मिळते. आता लक्षात घ्या की क्वाइनचे
खोटारड्याचे वाक्य हे “असत्य विधान मिळते” याचे नव्हे,
तर “उद्धृत रूप आधी जोडल्यास असत्य विधान मिळते” याचे
विकर्णरूप आहे. म्हणजेच खोटारड्याचे वाक्य मिळवण्यासाठी
ज्याचे विकर्णीकरण करतो, त्या गुणधर्मातच विकर्णीकरण
समाविष्ट असते!{}

अंकगणिताच्या भाषेत, एक मुक्त चल असलेल्या
सूत्राची गोडेल संख्या काढून, त्या चलाच्या जागी
त्या संख्येचा प्रमाण संख्यांक बसवून त्याचे उद्धृत रूप
बनवतो. $\mathsf{E}(x)$ चे विकर्णरूप हे $\mathsf{E}(\num{n})$ आहे,
जिथे $n = \Gn{\mathsf{E}(x)}$. (यापुढे
$\num{\Gn{\mathsf{E}(x)}}$ याचे संक्षिप्त रूप
$\gn{\mathsf{E}(x)}$ असे लिहू.) म्हणून $\psi(x)$ हा
“असत्य आहे” हा गुणधर्म असेल, तर “स्वतःचे उद्धृत रूप
आधी जोडल्यास असत्य विधान मिळते” याचा अर्थ
“स्वतःच्या विकर्णरूपाच्या गोडेल संख्येला लागू केल्यावर
असत्य विधान मिळते” असा होईल. सूत्राची गोडेल
संख्या~$n$ दिल्यावर त्याच्या विकर्णरूपाची गोडेल संख्या
काढणाऱ्या $\fn{diag}(n)$ या फलनासाठी
$\Obj{diag}$ असे चिन्ह आपल्याकडे असेल, तर
$\mathsf{E}(x)$ हे $\psi(\Obj{diag}(x))$ असे लिहिता येईल.
क्वाइनच्या खोटारड्याच्या वाक्याचे रूप मग याचे
विकर्णीकरण, म्हणजे $\mathsf{E}(\gn{\mathsf{E}(x)})$ किंवा
$\psi(\Obj{diag}(\gn{\psi(\Obj{diag}(x))}))$, असे असेल.
अर्थात, आता $\psi(x)$ हा इतर कोणताही गुणधर्म असला,
तरी हीच रचना चालेल. अपूर्णता प्रमेयासाठी $\psi(x)$
हा “$x$ हे~$\Th{T}$ मध्ये निष्पन्न करता येण्याजोगे नाही” हा
गुणधर्म घेऊ. तेव्हा $\mathsf{E}(x)$ चे वर्णन असे होईल:
“स्वतःच्या विकर्णरूपाच्या गोडेल संख्येला लागू केल्यावर
$\Th{T}$ मध्ये निष्पन्न करता येण्याजोगे नसलेले
वाक्य मिळते.”

$\Th{T}$ मध्ये ही रचना आकारिक रीतीने मांडण्यासाठी
$\fn{diag}$ चे आकारिकीकरण करावे लागेल.
$\fn{diag}(n)$ हे संगणनक्षम, किंबहुना आदिम पुनरावर्ती
फलन आहे: $n$ ही~$\mathsf{E}(x)$ या सूत्राची गोडेल
संख्या असेल, तर $\fn{diag}(n)$ हे
$\mathsf{E}(\gn{\mathsf{E}(x)})$ ची गोडेल संख्या देते. अन्य आदानांसाठी शून्य मूल्य निवडून
हे फलन सर्व नैसर्गिक संख्यांवर परिभाषित करतो.
(आठवा, $\gn{\mathsf{E}(x)}$ हा~$\mathsf{E}(x)$ च्या गोडेल
संख्येचा प्रमाण संख्यांक, म्हणजे
$\num{\Gn{\mathsf{E}(x)}}$, आहे.) $\Obj{diag}$ हे
$\fn{diag}$ चे निरूपण करणारे $\Th{T}$ मधील
फलनचिन्ह असेल, तर $\varphi$ म्हणून
$\psi(\Obj{diag}(\gn{\psi(\Obj{diag}(x))}))$
हे सूत्र घेता येईल. लक्षात घ्या:
\begin{align*}
\fn{diag}(\Gn{\psi(\Obj{diag}(x))}) & =
\Gn{\psi(\Obj{diag}(\gn{\psi(\Obj{diag}(x))}))} \\
& = \Gn{\varphi}.
\end{align*}
$\Th{T}$ पुढील विधान निष्पन्न करू शकते असे धरल्यास,
\[\Obj{diag}(\gn{\psi(\Obj{diag}(x))}) = \gn{\varphi},
\]
त्याला $\psi(\Obj{diag}(\gn{\psi(\Obj{diag}(x))}))
\liff \psi(\gn{\varphi})$ हेही निष्पन्न करता येते.
पण डावी बाजू व्याख्येनुसार~$\varphi$ आहे.

अर्थात, सामान्यतः $\Obj{diag}$ हे
$\Th{T}$ मधील फलनचिन्ह नसेल; आणि
$\Th{Q}$ मधील तर नक्कीच नाही. पण
$\fn{diag}$ संगणनक्षम असल्यामुळे ते
$\Th{Q}$ मध्ये एखाद्या
$\theta_{\fn{diag}}(x,y)$ सूत्राने
\emph{निरूपित} करता येते. म्हणून
$\psi(\Obj{diag}(x))$ असे लिहिण्याऐवजी
$\lexists[y][(\theta_{\fn{diag}}(x,y) \land \psi(y))]$
असे लिहू शकतो. याखेरीज वर रेखाटलेली सिद्धता तशीच
चालते; किंबहुना ती आधीच~$\Th{Q}$ मध्ये चालते.
\end{explain}

\begin{lem}
\label{inc:inp:fix:lem:fixed-point} $\psi(x)$ हे~$x$ हे एकमेव मुक्त चल
असलेले कोणतेही सूत्र असू द्या. मग असे वाक्य~$\varphi$
असते की $\Th{Q} \Proves
\varphi \liff \psi(\gn{\varphi})$.
\end{lem}


\begin{proof}
$\psi(x)$ दिल्यावर $\mathsf{E}(x)$ हे
$\lexists[y][(\theta_{\fn{diag}}(x,y) \land \psi(y))]$
सूत्र मानू, आणि $\varphi$ हे त्याचे विकर्णरूप,
म्हणजे $\mathsf{E}(\gn{\mathsf{E}(x)})$, मानू.

$\theta_{\fn{diag}}$ हे $\fn{diag}$ चे निरूपण
करते आणि $\fn{diag}(\Gn{\mathsf{E}(x)}) = \Gn{\varphi}$
असल्यामुळे $\Th{Q}$ पुढील विधाने
निष्पन्न करू शकते:
\begin{align}
  & \theta_{\fn{diag}}(\gn{\mathsf{E}(x)}, \gn{\varphi}) \label{inc:inp:fix:repdiag1} \\
  & \lforall[y][(\theta_{\fn{diag}}(\gn{\mathsf{E}(x)},y) \lif
  \eq[y][\gn{\varphi}])]. \label{inc:inp:fix:repdiag2}
\end{align}
आता $\Th{Q} \Proves \varphi \liff \psi(\gn{\varphi})$
हे दाखवू. केवळ तर्कशास्त्र आणि $\Th{Q}$
मध्ये निष्पन्न करता येण्याजोगे असलेली तथ्ये वापरून
अनौपचारिक युक्तिवाद करू.

प्रथम~$\varphi$, म्हणजे $\mathsf{E}(\gn{\mathsf{E}(x)})$, खरे
मानू. $\mathsf{E}(x)$ च्या व्याख्येकडे परत पाहिल्यास,
$\mathsf{E}(\gn{\mathsf{E}(x)})$ हेच पुढील विधान आहे:
\[\lexists[y][(\theta_{\fn{diag}}(\gn{\mathsf{E}(x)},y) \land \psi(y))].
\]
अशा~$y$ चा विचार करा. आता
$\theta_{\fn{diag}}(\gn{\mathsf{E}(x)},y)$ असल्यामुळे
\ref{inc:inp:fix:repdiag2} वरून $y = \gn{\varphi}$ मिळते.
म्हणून $\psi(y)$ वरून~$\psi(\gn{\varphi})$ मिळते.

आता $\psi(\gn{\varphi})$ मानू. \ref{inc:inp:fix:repdiag1}
वरून पुढील मिळते:
\begin{align*}
& \theta_{\fn{diag}}(\gn{\mathsf{E}(x)}, \gn{\varphi}) \land \psi(\gn{\varphi}).
\intertext{त्यावरून पुढील निष्पन्न होते:}
& \lexists[y][(\theta_{\fn{diag}}(\gn{\mathsf{E}(x)},y) \land \psi(y))].
\end{align*}
पण हेच $\mathsf{E}(\gn{\mathsf{E}(x)})$, म्हणजेच~$\varphi$, आहे.
\end{proof}

\begin{digress}
या सिद्धतेची संगणनक्षमता सिद्धांतातील स्थिरबिंदू
उपपत्तीच्या सिद्धतेशी तुलना करा. येथे स्वतःच्याच
साहाय्याने \emph{विधान} परिभाषित करायचे आहे,
तर तिथे स्वतःच्याच साहाय्याने \emph{फलन}
परिभाषित करायचे होते. हा फरक सोडला, तर
मूळ कल्पना तीच आहे.
\end{digress}

\begin{prob}
  सूत्र~$\varphi(x)$ ला \emph{सत्याची व्याख्या}
  असे म्हणतात, जर सर्व वाक्ये~$\psi$
  साठी $\Th{Q}
  \Proves \psi \liff \varphi(\gn{\psi})$ असेल.
  स्थिरबिंदू उपपत्ती वापरून कोणतेही
  सूत्र सत्याची व्याख्या असू शकत नाही
  हे दाखवा.
\end{prob}












\section{पहिले अपूर्णता प्रमेय}\label{inc:inp:1in:sec}

आता गोडेलने पहिल्या अपूर्णता प्रमेयासाठी दिलेल्या मूळ
सिद्धतेचे वर्णन करता येईल. अंकगणिताच्या भाषेचा विस्तार
असलेल्या भाषेतील $\Th{T}$ ही कोणतीही संगणनक्षमपणे
स्वयंसिद्धकीकृत उपपत्ती असू द्या, आणि $\Th{T}$ मध्ये
$\Th{Q}$ ची स्वयंसिद्धके समाविष्ट असू द्या. यावरून
विशेषतः $\Th{T}$ संगणनक्षम फलने आणि संबंध
निरूपित करते.

सूत्रे आणि सिद्धता संख्या म्हणून योग्य रीतीने सांकेतिक
केल्यास $\Prf[\Th{T}](x,y)$ हा संबंध संगणनक्षम
असतो, असा आपण युक्तिवाद केला आहे. येथे
$\Prf[\Th{T}](x,y)$ खरा असतो, तर आणि तरच
$x$ ही~$\Th{T}$ मध्ये गोडेल संख्या~$y$
असलेल्या सूत्राच्या निष्पत्ती
ची गोडेल संख्या असते. गोडेलच्या मनात असलेल्या
विशिष्ट उपपत्तीसाठी हा संबंध आदिम पुनरावर्ती आहे
हेही त्याने आपल्या शोधपत्रातील 45 फलने आणि
संबंधांची यादी वापरून दाखवले. त्या यादीतील
45 वा संबंध, $x B y$, हा त्याने निवडलेल्या
$\Th{T}$ साठीचा $\Prf[\Th{T}](x,y)$
संबंधच आहे. गोडेल आपल्या शोधपत्रात जिथे
“पुनरावर्ती” हा शब्द वापरतो, तिथे आज आपण
“आदिम पुनरावर्ती” ही संज्ञा वापरू, हे लक्षात ठेवा.

$\Prf[\Th{T}](x,y)$ संगणनक्षम असल्यामुळे
तो $\Th{T}$ मध्ये निरूपणीय आहे. त्याचे
निरूपण करणाऱ्या सूत्रासाठी
$\OPrf[\Th{T}](x,y)$ हे चिन्ह वापरू.
$\OProv[\Th{T}](y)$ हे सूत्र
$\lexists[x][\OPrf[\Th{T}](x,y)]$
असे मानू. हे गोडेलच्या यादीतील 46 व्या
संबंधाचे, $ \fn{Bew}(y)$ चे, वर्णन करते.
गोडेलच्या नोंदीनुसार, हाच एक संबंध असा
आहे की तो “पुनरावर्ती आहे असे म्हणता येत
नाही.” त्याचा संभाव्य अर्थ असा: व्याख्येवरून
तो संगणनक्षम आहे हे स्पष्ट होत नाही; आणि
नंतरच्या प्रगतीतून, उपपत्ती क्यूचा सुसंगत विस्तार असेल,
तर हा संबंध संगणनक्षम नसतो हेही सिद्ध झाले.

$\Th{T}$ ही~$\Th{Q}$ समाविष्ट करणारी
स्वयंसिद्धकीकरणीय उपपत्ती असू द्या. मग
$\Prf[\Th{T}](x, y)$ निर्णेय असतो आणि
म्हणून $\Th{Q}$ मध्ये
सूत्र~$\OPrf[\Th{T}](x, y)$
याने निरूपणीय असतो. वर वर्णन केल्याप्रमाणे
$\OProv[\Th{T}](y)$ हे सूत्र मानू.
स्थिरबिंदू उपपत्तीवरून असे सूत्र
$\mathsf{G}_\Th{T}$ मिळते की $\Th{Q}$
(आणि म्हणून $\Th{T}$ सुद्धा) पुढील
निष्पन्न करते:
\begin{equation}
\label{inc:inp:1in:eqn:qpf}
\mathsf{G}_\Th{T} \liff \lnot \OProv[\Th{T}](\gn{\mathsf{G}_\Th{T}}).
\end{equation}
लक्षात घ्या की $\mathsf{G}_\Th{T}$ चे
आशयतः विधान असे आहे: “$\mathsf{G}_\Th{T}$
हे~$\Th{T}$ मध्ये निष्पन्न करता येण्याजोगे नाही.”

\begin{lem}\label{inc:inp:1in:lem:cons-G-unprov}
$\Th{T}$ ही~$\Th{Q}$ चा विस्तार असलेली
सुसंगत आणि स्वयंसिद्धकीकरणीय उपपत्ती
असेल, तर $\Th{T} \Proves/ \mathsf{G}_\Th{T}$.
\end{lem}

\begin{proof}
$\Th{T}$ हे $\mathsf{G}_\Th{T}$ ला
निष्पन्न करते असे मानू. मग
निष्पत्ती \emph{अस्तित्वात आहे}.
म्हणून एखाद्या संख्या $m$ साठी
$\Prf[\Th{T}](m,\Gn{\mathsf{G}_\Th{T}})$
हा संबंध खरा असतो. पण मग
$\Th{Q}$ हे
$\OPrf[\Th{T}](\num m, \gn{\mathsf{G}_\Th{T}})$
वाक्य निष्पन्न करते.
त्यामुळे $\Th{Q}$ हे
$\lexists[x][\OPrf[\Th{T}](x,\gn{\mathsf{G}_\Th{T}})]$
सुद्धा निष्पन्न करते; व्याख्येनुसार
हेच $\OProv[\Th{T}](\gn{\mathsf{G}_\Th{T}})$
आहे. \ref{inc:inp:1in:eqn:qpf} वरून
$\Th{Q}$ हे $\lnot \mathsf{G}_\Th{T}$
निष्पन्न करते. $\Th{T}$ हा
$\Th{Q}$ चा विस्तार असल्यामुळे
$\Th{T}$ सुद्धा तसेच करते.
म्हणजे $\Th{T}$ हे $\mathsf{G}_\Th{T}$
निष्पन्न करत असेल, तर
$\lnot \mathsf{G}_\Th{T}$ सुद्धा
निष्पन्न करते, असे आपण दाखवले.
त्यामुळे ते विसंगत ठरेल.
\end{proof}

\begin{defn}
\label{inc:inp:1in:thm:oconsis-q}
उपपत्ती $\Th{T}$ ला
\emph{$\omega$-सुसंगत} म्हणतात, जर
पुढील अट पूर्ण होत असेल: कोणतेही
वाक्य $\lexists[x][\varphi(x)]$ घेऊ.
$\Th{T}$ हे $\lnot \varphi(\num 0)$,
$\lnot \varphi(\num 1)$,
$\lnot \varphi(\num 2)$, \dots
निष्पन्न करत असेल, तर
$\Th{T}$ हे $\lexists[x][\varphi(x)]$ सिद्ध करत नाही.
\end{defn}

प्रत्येक $\omega$-सुसंगत उपपत्ती
सुसंगतसुद्धा असते, हे लक्षात घ्या.
कारण $\Th{T}$ विसंगत असेल, तर
प्रत्येक~$\varphi$ साठी $\Th{T} \Proves \varphi$
असते. विशेषतः $\Th{T}$ विसंगत
असेल, तर प्रत्येक~$n$ साठी
$\lnot \varphi(\num n)$ हे ते
निष्पन्न करते आणि
$\lexists[x][\varphi(x)]$ हेही
निष्पन्न करते.
म्हणून विसंगत $\Th{T}$ हे
$\omega$-विसंगत असते. याच्या
प्रतिवर्ती विधानावरून
$\omega$-सुसंगत $\Th{T}$
सुसंगत असलेच पाहिजे.

\begin{lem}\label{inc:inp:1in:lem:omega-cons-G-unref}
$\Th{T}$ ही~$\Th{Q}$ चा विस्तार
असलेली $\omega$-सुसंगत,
स्वयंसिद्धकीकरणीय उपपत्ती
असेल, तर
$\Th{T} \Proves/ \lnot \mathsf{G}_\Th{T}$.
\end{lem}

\begin{proof}
$\Th{T}$ हे $\lnot \mathsf{G}_\Th{T}$
निष्पन्न करत असेल, तर ते
$\omega$-विसंगत असते, हे दाखवू.
$\Th{T}$ हे $\lnot \mathsf{G}_\Th{T}$
निष्पन्न करते असे मानू.
$\Th{T}$ विसंगत असेल, तर ते
$\omega$-विसंगत आहेच.
नसेल, तर $\Th{T}$ सुसंगत आहे.
म्हणून \ref{inc:inp:1in:lem:cons-G-unprov}
वरून ते $\mathsf{G}_\Th{T}$ ला
निष्पन्न करत नाही. $\Th{T}$
मध्ये $\mathsf{G}_\Th{T}$ ची
निष्पत्ती नसल्यामुळे
$\Th{Q}$ पुढील विधाने
निष्पन्न करते:
\[\lnot \OPrf[\Th{T}](\num 0, \gn{\mathsf{G}_\Th{T}}), \lnot \OPrf[\Th{T}](\num 1,
\gn{\mathsf{G}_\Th{T}}), \lnot \OPrf[\Th{T}](\num 2, \gn{\mathsf{G}_\Th{T}}), \dots
\]
त्यामुळे $\Th{T}$ सुद्धा ती
सिद्ध करते. दुसरीकडे
\ref{inc:inp:1in:eqn:qpf} वरून
$\lnot \mathsf{G}_\Th{T}$ हे
$\lexists[x][\OPrf[\Th{T}](x,\gn{\mathsf{G}_\Th{T}})]$
याच्याशी सममूल्य असते.
म्हणून $\Th{T}$ हे
$\omega$-विसंगत आहे.
\end{proof}

\begin{prob}
  प्रत्येक $\omega$-सुसंगत उपपत्ती
  सुसंगत असते. याचे उलट खरे नसते
  हे दाखवा: सुसंगत पण
  $\omega$-विसंगत उपपत्ती असतात.
  यासाठी $\Th{Q} \cup
  \{\lnot \mathsf{G}_\Th{Q}\}$ हा संच
  सुसंगत पण $\omega$-विसंगत
  आहे असे दाखवा.
\end{prob}

\begin{thm}
\label{inc:inp:1in:thm:first-incompleteness}
$\Th{T}$ ही~$\Th{Q}$ चा
$\omega$-सुसंगत,
स्वयंसिद्धकीकरणीय विस्तार असलेली
कोणतीही उपपत्ती असू द्या.
मग $\Th{T}$ संपूर्ण नाही.
\end{thm}

\begin{proof}
  $\Th{T}$ हे $\omega$-सुसंगत
  असेल, तर ते सुसंगत असते.
  म्हणून \ref{inc:inp:1in:lem:cons-G-unprov}
  वरून $\Th{T}
  \Proves/ \mathsf{G}_\Th{T}$.
  \ref{inc:inp:1in:lem:omega-cons-G-unref}
  वरून $\Th{T} \Proves/
  \lnot \mathsf{G}_\Th{T}$. म्हणजे
  $\Th{T}$ अपूर्ण आहे: ते
  $\mathsf{G}_\Th{T}$ किंवा
  $\lnot \mathsf{G}_\Th{T}$ यांपैकी
  एकही विधान निष्पन्न करत नाही.
\end{proof}












\section{रोसरचे प्रमेय}\label{inc:inp:ros:sec}

गोडेलची सिद्धता बदलून अधिक सबळ निकाल मिळेल का,
म्हणजे “$\omega$-सुसंगत” ऐवजी फक्त
“सुसंगत” असे गृहीत धरता येईल का?
रोसरने शोधलेल्या युक्तीमुळे याचे उत्तर
“होय” असे आहे. रोसरची युक्ती म्हणजे
$\OProv[\Th{T}](y)$ च्या जागी
“सुधारित” निष्पन्नता विधेय
$\ORProv_T(y)$ वापरणे.

\begin{thm}
\label{inc:inp:ros:thm:rosser}
$\Th{T}$ ही~$\Th{Q}$ चा विस्तार
असलेली कोणतीही सुसंगत आणि
स्वयंसिद्धकीकरणीय उपपत्ती असू द्या.
मग $\Th{T}$ संपूर्ण नसते.
\end{thm}

\begin{proof}
आठवा, $\OProv[\Th{T}](y)$ ची व्याख्या
$\lexists[x][\OPrf[\Th{T}](x,
  y)]$ अशी आहे. येथे
$\OPrf[\Th{T}](x, y)$ हे अशा निर्णेय
संबंधाचे निरूपण करते की $x$ ही
गोडेल संख्या~$y$ असलेल्या
वाक्याच्या निष्पत्ती
ची गोडेल संख्या असते, तर आणि तरच तो
संबंध खरा असतो. $x$ आणि~$y$
यांच्यातील पुढील संबंधही निर्णेय
असतो: $x$ ही गोडेल
संख्या~$y$ असलेल्या वाक्याच्या
\emph{खंडनाची} गोडेल संख्या असते.
$\fn{not}(x)$ हे पुढील कार्य
करणारे आदिम पुनरावर्ती फलन
असू द्या: $x$ हा $\varphi$ या
सूत्राचा संकेतांक असेल, तर
$\fn{not}(x)$ हा $\lnot
\varphi$ चा संकेतांक असतो. सूत्रसंकेतांक नसलेल्या
आदानासाठी शून्य परत देऊन हे फलन सर्वत्र परिभाषित करतो. मग
$\Refut[\Th{T}](x, y)$
खरे असते, तर आणि तरच
$\Prf[\Th{T}](x, \fn{not}(y))$
खरे असते. त्याचे निरूपण
$\ORefut[\Th{T}](x, y)$
करू द्या. मग
$\Th{T} \Proves \lnot \varphi$
असेल आणि $\delta$ ही
संबंधित निष्पत्ती असेल,
तर $\Th{Q} \Proves
\ORefut[\Th{T}](\gn{\delta}, \gn{\varphi})$.
आता $\ORProv[\Th{T}](y)$ ची
व्याख्या अशी करू:
\[\lexists[x][(\OPrf[\Th{T}](x,y) \land \lforall[z][(z < x \lif \lnot
  \ORefut[\Th{T}](z,y))])].
\]
स्थूलपणे, $\ORProv[\Th{T}](y)$
असे सांगते की “$y$ ची
$\Th{T}$ मध्ये सिद्धता आहे,
आणि कमी संकेतांक असलेले
$y$ चे खंडन नाही.”
$\Th{T}$ सुसंगत असेल,
तर $\ORProv[\Th{T}](y)$ आणि
$\OProv[\Th{T}](y)$ ही दोन्ही
विधाने त्याच संख्यांसाठी
सत्य असतात. पण $\Th{T}$
मधील \emph{सिद्धतायोग्यतेच्या}
दृष्टीने (आता सत्य आणि
सिद्धतायोग्यता यांत फरक आहे
हे आपल्याला माहीत आहे)
त्यांचे गुणधर्म भिन्न असतात.
$\Th{T}$ \emph{वि}संगत
असेल, तर ती दोन्ही विधाने
त्याच संख्यांसाठी सत्य
\emph{नाहीत}!{}
($\ORProv[\Th{T}](y)$
याचा अर्थ अनेकदा “$y$ रोसर
अर्थाने सिद्धतायोग्य आहे”
असा घेतला जातो. पण आत्ताच
पाहिल्याप्रमाणे रोसर
सिद्धतायोग्यता हा
सिद्धतायोग्यतेचा काही
विशेष प्रकार नाही:
विसंगत उपपत्तीत काही
वाक्ये सिद्धतायोग्य
असतात, पण रोसर अर्थाने
सिद्धतायोग्य नसतात.
म्हणून हा शब्द गोंधळात
टाकू शकतो. तो टाळण्यासाठी
“$y$ शमोव्हेबल आहे”
असेही म्हणता येईल.)

स्थिरबिंदू उपपत्तीवरून
असे सूत्र $\mathsf{R}_\Th{T}$
असते की
\begin{equation}
  \Th{Q} \Proves \mathsf{R}_\Th{T} \liff \lnot \ORProv[\Th{T}](\gn{\mathsf{R}_\Th{T}}).
  \label{inc:inp:ros:RT}
\end{equation}
\ref{inc:inp:1in:thm:first-incompleteness}
च्या सिद्धतेपेक्षा इथे
आपला दावा अधिक सबळ आहे:
$\Th{T}$ सुसंगत असेल,
तर $\Th{T}$ हे
$\mathsf{R}_\Th{T}$ ला निष्पन्न
करत नाही, आणि $\Th{T}$
$\lnot \mathsf{R}_\Th{T}$ लाही
निष्पन्न करत नाही.
(म्हणजे $\omega$-सुसंगततेचे
गृहीतक लागत नाही.)

प्रथम $\Th{T} \Proves/
\mathsf{R}_{\Th{T}}$ हे दाखवू.
उलट ते सिद्ध होत असेल,
तर~$T$ पासून
$\mathsf{R}_{\Th{T}}$ ची
निष्पत्ती आहे;
तिची गोडेल संख्या $n$
मानू. $\OPrf[\Th{T}]$
हे $\Prf[\Th{T}]$ चे
$\Th{Q}$ मध्ये निरूपण
करत असल्यामुळे
$\Th{Q} \Proves
\OPrf[\Th{T}](\num{n}, \gn{\mathsf{R}_{\Th{T}}})$.
शिवाय प्रत्येक $k < n$
साठी, $\Th{T}$ सुसंगत
असल्यामुळे, $k$ ही
$\lnot \mathsf{R}_{\Th{T}}$
च्या निष्पत्तीची
गोडेल संख्या नसते.
म्हणून प्रत्येक $k < n$
साठी $\Th{Q} \Proves \lnot
\ORefut[\Th{T}](\num{k}, \gn{\mathsf{R}_{\Th{T}}})$.
\ref{inc:req:min:lem:less-nsucc}
वरून $\Th{Q} \Proves
\lforall[z][(z < \num{n} \lif \lnot \ORefut[\Th{T}](z,
  \gn{\mathsf{R}_{\Th{T}}}))]$.
त्यामुळे
\[\Th{Q} \Proves \lexists[x][(\OPrf[\Th{T}](x,\gn{\mathsf{R}_{\Th{T}}}) \land \lforall[z][(z
    < x \lif \lnot \ORefut[\Th{T}](z,\gn{\mathsf{R}_{\Th{T}}}))])],
\]
पण हेच
$\ORProv[\Th{T}](\gn{\mathsf{R}_{\Th{T}}})$
आहे. \ref{inc:inp:ros:RT} वरून
$\Th{Q}
\Proves \lnot \mathsf{R}_{\Th{T}}$.
$\Th{T}$ हा $\Th{Q}$ चा
विस्तार असल्यामुळे
$\Th{T}
\Proves \lnot \mathsf{R}_{\Th{T}}$
सुद्धा मिळते. पण आपण
$\Th{T} \Proves \mathsf{R}_{\Th{T}}$
असे गृहीत धरले होते.
त्यामुळे $\Th{T}$ विसंगत
ठरेल; हे प्रमेयाच्या
गृहीतकाच्या विरुद्ध आहे.

आता $\Th{T} \Proves/
\lnot \mathsf{R}_{\Th{T}}$ हे
दाखवू. उलट ते सिद्ध होत
असेल आणि
$\lnot \mathsf{R}_{\Th{T}}$
च्या निष्पत्तीची
गोडेल संख्या $n$ असेल
असे मानू. मग
$\Refut[\Th{T}](n, \Gn{\mathsf{R}_{\Th{T}}})$
खरे असते. $\ORefut[\Th{T}]$
हे $\Refut[\Th{T}]$ चे
$\Th{Q}$ मध्ये निरूपण
करत असल्यामुळे
$\Th{Q} \Proves
\ORefut[\Th{T}](\num{n}, \gn{\mathsf{R}_{\Th{T}}})$.
मग $\Th{T}$ हे
$\mathsf{R}_{\Th{T}}$ सुद्धा
निष्पन्न करेल आणि
त्यामुळे विसंगत ठरेल,
हे पुन्हा दाखवू. कारण
\begin{align*}
\Th{Q} & \Proves \mathsf{R}_{\Th{T}} \liff \lnot \ORProv[\Th{T}](\gn{\mathsf{R}_{\Th{T}}}),
\intertext{आणि $\Th{T}$ हा~$\Th{Q}$ चा विस्तार असल्यामुळे पुढील दाखवणे पुरेसे आहे:}
\Th{Q} & \Proves \lnot \ORProv[\Th{T}](\gn{\mathsf{R}_{\Th{T}}}).
\end{align*}
वाक्य $\lnot
\ORProv[\Th{T}](\gn{\mathsf{R}_{\Th{T}}})$,
म्हणजे
\begin{align*}
  \lnot & \lexists[x][(\OPrf[\Th{T}](x,\gn{\mathsf{R}_{\Th{T}}}) \land
    \lforall[z][(z < x \lif \lnot \ORefut[\Th{T}](z,\gn{\mathsf{R}_{\Th{T}}}))])],
  \intertext{हे तर्कशास्त्रीय दृष्ट्या पुढील विधानाशी सममूल्य आहे:}
  & \lforall[x][(\OPrf[\Th{T}](x,\gn{\mathsf{R}_{\Th{T}}}) \lif
    \lexists[z][(z < x \land \ORefut[\Th{T}](z,\gn{\mathsf{R}_{\Th{T}}}))])].
\end{align*}
$\Th{Q}$ काय निष्पन्न
करते याविषयीची तथ्ये वापरून
अनौपचारिक युक्तिवाद करू.
$x$ स्वेच्छेने निवडलेला
आणि $\OPrf[\Th{T}](x,
\gn{\mathsf{R}_{\Th{T}}})$ खरा
आहे असे मानू. आपल्याला
आधीच माहीत आहे की
$\Th{T} \Proves/ \mathsf{R}_{\Th{T}}$.
म्हणून प्रत्येक $k$
साठी $\Th{Q} \Proves \lnot
\OPrf[\Th{T}](\num{k}, \gn{\mathsf{R}_{\Th{T}}})$.
त्यामुळे प्रत्येक $k$
साठी $\eq/[x][\num{k}]$
मिळते. विशेषतः (a)
$\eq/[x][\num{n}]$
मिळते. तसेच
$\lnot(\eq[x][\num{0}]
\lor \eq[x][\num{1}] \lor \dots
\lor \eq[x][\num{n-1}])$
मिळते. म्हणून
\ref{inc:req:min:lem:less-nsucc}
वरून (b) $\lnot(x < \num{n})$.
\ref{inc:req:min:lem:trichotomy}
वरून $\num{n} < x$.
$\Th{Q} \Proves
\ORefut[\Th{T}](\num{n}, \gn{\mathsf{R}_{\Th{T}}})$
असल्यामुळे
$\num{n} < x \land
\ORefut[\Th{T}](\num{n}, \gn{\mathsf{R}_{\Th{T}}})$
मिळते, आणि त्यावरून
$\lexists[z][(z < x
\land \ORefut[\Th{T}](z,\gn{\mathsf{R}_{\Th{T}}}))]$
मिळते. $x$ स्वेच्छेने
निवडलेला असल्यामुळे
अपेक्षेप्रमाणे पुढील मिळते:
\[\lforall[x][(\OPrf[\Th{T}](x,\gn{\mathsf{R}_{\Th{T}}}) \lif
  \lexists[z][(z < x \land \ORefut[\Th{T}](z,\gn{\mathsf{R}_{\Th{T}}}))])].
\]
\end{proof}

\begin{prob}
नैसर्गिक संख्यांचे दोन संच
$A$ आणि $B$ यांना
\emph{संगणनक्षमपणे अविलगनीय}
म्हणतात, जर $A \subseteq X$
आणि $B \subseteq \Complement{X}$
असेल असा कोणताही
निर्णेय संच~$X$
नसेल ($\Complement{X}$
म्हणजे~$X$ चा पूरक
$\Nat \setminus X$).
$\Th{T}$ ही~$\Th{Q}$
ची सुसंगत,
स्वयंसिद्धकीकरणीय विस्तार-उपपत्ती
असू द्या. $\Th{T}$
मध्ये सिद्ध होणाऱ्या
वाक्यांच्या गोडेल
संख्यांचा संच $A$,
आणि $\Th{T}$ मध्ये
खंडित होणाऱ्या वाक्यांच्या
गोडेल संख्यांचा संच $B$
असेल असे मानू. $A$ आणि
$B$ हे संगणनक्षमपणे
अविलगनीय आहेत हे सिद्ध करा.
\end{prob}












\section{गोडेलच्या मूळ शोधपत्राशी तुलना}\label{inc:inp:gop:sec}

गोडेलच्या 1931 मधील शोधपत्राला थोडा वेळ
देणे उपयुक्त ठरेल. त्याच्या प्रस्तावनेत आपण
नुकत्याच पाहिलेल्या कल्पनांची रूपरेषा आहे.
शोधपत्र नुसते चाळले तरी प्रत्येक टप्प्यावर
काय चालले आहे हे सहज दिसते: प्रथम गोडेल
आकारिक पद्धत $P$ (विन्यास, स्वयंसिद्धके,
सिद्धतांचे नियम) वर्णन करतो; मग आदिम
पुनरावर्ती फलने आणि संबंध परिभाषित करतो;
त्यानंतर $x B y$ हा संबंध आदिम पुनरावर्ती
आहे हे दाखवतो आणि आदिम पुनरावर्ती फलने
व संबंध $\Th{P}$ मध्ये निरूपित होतात
असा युक्तिवाद करतो. मग वर दिल्याप्रमाणे
अपूर्णता प्रमेय सिद्ध करतो. तिसऱ्या विभागात,
सिद्ध करता न येणारे विधान अंकगणिताच्या
भाषेतील वाक्य म्हणून घेता येते हे दाखवतो.
यातूनच $\beta$-उपपत्ती उद्भवली; पुनरावर्ती
फलने $\Th{Q}$ मध्ये निरूपणीय आहेत हे
दाखवताना अनुक्रम हाताळण्यासाठी आपणही ती
वापरली. गोडेल एवढेच पुरेसे ठरणाऱ्या
किमान स्वयंसिद्धकसंचाची वेगळी ओळख करून
देत नाही, पण आता $\Th{Q}$ ते काम करते
हे आपल्याला माहीत आहे. अखेरीस चौथ्या
विभागात तो दुसऱ्या अपूर्णता प्रमेयाच्या
सिद्धतेची रूपरेषा देतो.












\section[PA साठीच्या निष्पन्नता अटी]{$\Th{PA}$ साठीच्या निष्पन्नता अटी}\label{inc:inp:prc:sec}

पेआनो अंकगणित, म्हणजे $\Th{PA}$, ही
सर्व सूत्रे साठी विगमनाची स्वयंसिद्धके
जोडून मिळणारी $\Th{Q}$ ची विस्तार-उपपत्ती
आहे. म्हणजे प्रत्येक सूत्र~$\varphi$
साठी $\Th{Q}$ मध्ये पुढील रूपाचे
स्वयंसिद्धक जोडतात:
\[(\varphi(0) \land \lforall[x][(\varphi(x) \lif \varphi(x'))]) \lif \lforall[x][\varphi(x)]
\]
या रूपबंधात इतर मुक्त चल असतील, तर आरंभी त्यांच्यावर
सार्विक संख्यापक जोडून वाक्यरूपाचे स्वयंसिद्धक घेतो.
हे खरे तर एक \emph{रूपबंध} आहे, म्हणजे
अनंत स्वयंसिद्धके आहेत, हे लक्षात घ्या
(आणि $\Th{PA}$ चे सांत
स्वयंसिद्धकसंचाने स्वयंसिद्धकीकरण करता येत
{\em नाही} असेही सिद्ध होते).
पण चिन्हमाला विगमनाच्या
स्वयंसिद्धकाचे उदाहरण आहे
की नाही हे प्रभावीपणे ठरवता
येत असल्यामुळे $\Th{PA}$
चा स्वयंसिद्धकसंच संगणनक्षम
आहे. $\Th{PA}$ ही~$\Th{Q}$
पेक्षा बरीच अधिक सबल
उपपत्ती आहे. उदाहरणार्थ,
नेहमीच्या पद्धतीने विगमन
वापरून बेरीज आणि गुणाकार
क्रमनिरपेक्ष आहेत हे
सहज सिद्ध करता येते.
खरे तर बहुतेक सांत
संख्यासिद्धांतिक आणि
संयोजनात्मक युक्तिवाद
$\Th{PA}$ मध्ये मांडता येतात.

$\Th{PA}$ ही संगणनक्षमपणे
स्वयंसिद्धकीकृत असल्यामुळे
$\Prf[\Th{PA}](x,y)$ हे
निष्पन्नता विधेय
संगणनक्षम आहे आणि म्हणून
$\Th{Q}$ मध्ये निरूपित
होते (म्हणूनच
$\Th{PA}$ मध्येही).
पूर्वीप्रमाणे, या संबंधाचे
निरूपण करणाऱ्या सूत्रासाठी
$\OPrf[\Th{PA}](x,y)$
वापरू. $\OProv[\Th{PA}](y)$
हे सूत्र
$\lexists[x][\OPrf[\Th{PA}](x,y)]$
असे मानू. त्याचा सहज
अर्थ असा: “$y$ हे
$\Th{PA}$ च्या स्वयंसिद्धकांपासून
निष्पन्न करता येण्याजोगे आहे.”
$\Th{Q}$ च्या स्वयंसिद्धकांपेक्षा
थोडे अधिक सामर्थ्य लागते,
कारण आपण वापरत असलेली
उपपत्ती या निष्पन्नता
विधेयाविषयीची काही मूलभूत
तथ्ये निष्पन्न करू शकली
पाहिजे. पुढील अटींसाठी प्रत्यक्ष आकारिक सिद्धतांच्या प्रमाण
संकेतांकीकरणावर आधारित विधेय निवडले आहे. केवळ निर्णेय
संबंधाचे निरूपण करणारे मनमानी सूत्र घेतल्याने या अटी
आपोआप मिळत नाहीत. नेमकी आवश्यक तथ्ये
पुढीलप्रमाणे:
\begin{enumerate}
\item[P1.] जर $\Th{PA} \Proves \varphi$, तर $\Th{PA} \Proves
  \OProv[\Th{PA}](\gn{\varphi})$.
\item[P2.] सर्व सूत्रे $\varphi$ आणि $\psi$ साठी,
  \[  \Th{PA} \Proves \OProv[\Th{PA}](\gn{\varphi \lif \psi}) \lif
  (\OProv[\Th{PA}](\gn{\varphi}) \lif \OProv[\Th{PA}](\gn{\psi})).
  \]
\item[P3.] प्रत्येक सूत्र~$\varphi$ साठी,
  \[  \Th{PA} \Proves \OProv[\Th{PA}](\gn{\varphi})
  \lif \OProv[\Th{PA}](\gn{\OProv[\Th{PA}](\gn{\varphi})}).
  \]
\end{enumerate}
हे तीन गुणधर्म खरे आहेत
याची खात्री करण्यासाठी
$\OProv[\Th{PA}](y)$
या सूत्राचे बारकाईने
वर्णन करून संबंधित आकारिक
निष्पत्त्यांचे वर्णन
$\Th{PA}$ च्या स्वयंसिद्धकांनी
करावे लागते. अटी (1) आणि~(2)
सोप्या आहेत; खरे काम
अट~(3) साठी लागते
(त्यासाठी नेमके काय
करावे लागेल याचा विचार करा
\dots). तपशील मांडणे
कंटाळवाणे आणि फारसे
रंजक नसल्यामुळे
$\Th{PA}$ मध्ये वरचे
तीन गुणधर्म आहेत हे
येथे आपण गृहीत धरू.
$\OProv[\Th{PA}](y)$
ची योग्य निवड पुढील
अटही पूर्ण करेल:
\begin{enumerate}
\item[P4.] जर $\Th{PA} \Proves \OProv[\Th{PA}](\gn{\varphi})$, तर
  $\Th{PA} \Proves \varphi$.
\end{enumerate}
पण हे तथ्य आपल्याला
लागणार नाही. ही शेवटची अट आधीच्या तीन आकारिक अटींचा
निष्कर्ष नाही; तिच्यासाठी पेआनो अंकगणिताचे मानक प्रतिमानातील
सत्य आणि प्रत्यक्ष सिद्धतासंकेतांकांचे प्रमाण विधेय वापरले आहे.

\begin{digress}
तपशील वगळण्याबाबत
गोडेलनेही आपल्यासारखाच
आळस केला. 1931 च्या
शोधपत्राच्या अखेरीस
त्याने दुसऱ्या अपूर्णता
प्रमेयाच्या सिद्धतेची
रूपरेषा दिली आणि
तपशील पुढील शोधपत्रात
देण्याचे आश्वासन दिले.
ते तपशील त्याने
नंतर प्रकाशित केले
नाहीत; युक्तिवाद समजणाऱ्या
सर्वांनाच तो पूर्ण करता
येईल असे वाटत असल्यामुळे
तपशील भरण्याची गरजही
त्याला पडली नाही.
\end{digress}












\section{दुसरे अपूर्णता प्रमेय}\label{inc:inp:2in:sec}

$\Th{PA}$ स्वतःची सुसंगतता सिद्ध करू
शकत नाही, हा दावा कसा व्यक्त करायचा?
$\Th{PA}$ विसंगत आहे असे म्हणणे
म्हणजे $\Th{PA} \Proves \eq[0][1]$
असे म्हणणे. सुसंगतता-विधानासाठी
असत्यता-स्थिराचा संकेतांक घेतला तरी
सुसंगततेचा आशय तोच राहतो; या
सिद्धतेत $\OCon[\Th{PA}]$ हे
वाक्य $\lnot
\OProv[\Th{PA}](\gn{\lfalse})$
मानू. मग पुढील प्रमेय अपेक्षित
निकाल देते:

\begin{thm}
\label{inc:inp:2in:thm:second-incompleteness}
$\Th{PA}$ सुसंगत असेल, तर
$\Th{PA}$ हे
$\OCon[\Th{PA}]$ निष्पन्न
करत नाही.
\end{thm}

हे प्रमेय $\OCon[\Th{PA}]$
च्या विशिष्ट निरूपणावर
(म्हणजे $\OProv[\Th{PA}](y)$
च्या विशिष्ट निरूपणावर)
अवलंबून आहे, हे लक्षात
घेणे महत्त्वाचे आहे.
आपण फक्त एवढेच वापरणार
आहोत की $\OProv[\Th{PA}](y)$
चे निरूपण तीन
निष्पन्नता अटी पूर्ण
करते. म्हणून असे गुणधर्म
असलेले निष्पन्नता
विधेय असलेल्या, क्यूचा विस्तार असलेल्या कोणत्याही
सुसंगत उपपत्तीला हे प्रमेय लागू
होते.

गोडेलच्या युक्तिवादाची
रूपरेषा वाचणे उपयुक्त आहे;
प्रमेय त्यातून नेमके
समारोपाप्रमाणे निघते.
युक्तिवाद असा आहे.
\ref{inc:inp:1in:thm:first-incompleteness}
च्या सिद्धतेत आपण तयार
केलेले गोडेल-वाक्य
$\mathsf{G}_\Th{PA}$ असू द्या.
“$\Th{PA}$ सुसंगत
असेल, तर $\Th{PA}$
हे $\mathsf{G}_\Th{PA}$ ला
निष्पन्न करत नाही”
हे आपण दाखवले आहे.
हे $\Th{PA}$
\emph{मध्येच} आकारिक
केल्यास पुढील विधानाची
सिद्धता मिळते:
\[\OCon[\Th{PA}] \lif \lnot \OProv[\Th{PA}](\gn{\mathsf{G}_\Th{PA}}).
\]
आता $\Th{PA}$ हे
$\OCon[\Th{PA}]$
निष्पन्न करते असे
मानू. मग ते $\lnot
\OProv[\Th{PA}](\gn{\mathsf{G}_\Th{PA}})$
सुद्धा निष्पन्न
करते. पण $\mathsf{G}_\Th{PA}$
हे गोडेल-वाक्य असल्यामुळे
हे $\mathsf{G}_\Th{PA}$ शी
सममूल्य आहे. म्हणून
$\Th{PA}$ हे
$\mathsf{G}_\Th{PA}$
निष्पन्न करते.

पण $\Th{PA}$ सुसंगत
असेल, तर ते
$\mathsf{G}_\Th{PA}$ ला
निष्पन्न करत
नाही, हे आपल्याला
माहीत आहे!{}
म्हणून $\Th{PA}$
सुसंगत असेल, तर
$\OCon[\Th{PA}]$
सुद्धा निष्पन्न
करू शकत नाही.

युक्तिवाद अधिक नेमका
करण्यासाठी $\mathsf{G}_\Th{PA}$
हे~$\Th{PA}$ चे
गोडेल-वाक्य मानू.
निष्पन्नता अटी
(P1)--(P3) वापरून
$\Th{PA}$ हे
$\OCon[\Th{PA}] \lif
\mathsf{G}_\Th{PA}$ निष्पन्न
करते हे दाखवू.
त्यावरून $\Th{PA}$
हे $\OCon[\Th{PA}]$
निष्पन्न करत नाही
हे मिळेल. पुढे
$\Th{PA}$ मधील
सिद्धतेची रूपरेषा आहे.
(सोयीसाठी
$\Th{PA}$ हे अधोलिखित
वगळू.)
\begin{align}
& \mathsf{G} \liff \lnot \OProv(\gn{\mathsf{G}}) \label{inc:inp:2in:G2-1}\\
& \qquad\text{$\mathsf{G}$ हे गोडेल-वाक्य आहे}\notag \\
& \mathsf{G} \lif \lnot \OProv(\gn{\mathsf{G}}) \label{inc:inp:2in:G2-2}\\
  & \qquad\text{\ref{inc:inp:2in:G2-1} वरून} \notag\\
& \mathsf{G} \lif
  (\OProv(\gn{\mathsf{G}}) \lif \lfalse) \label{inc:inp:2in:G2-3}\\
  & \qquad\text{\ref{inc:inp:2in:G2-2} व तर्कशास्त्रावरून}\notag\\
& \OProv(\gn{
    \mathsf{G} \lif
    (\OProv(\gn{\mathsf{G}}) \lif \lfalse)
  }) \label{inc:inp:2in:G2-4}\\
  & \qquad\text{\ref{inc:inp:2in:G2-3} आणि अट P1 वरून} \notag\\
& \OProv(\gn{\mathsf{G}}) \lif
  \OProv(\gn{
    (\OProv(\gn{\mathsf{G}}) \lif \lfalse)
    }) \label{inc:inp:2in:G2-5}\\
  & \qquad\text{\ref{inc:inp:2in:G2-4} आणि अट P2 वरून} \notag\\
& \OProv(\gn{\mathsf{G}}) \lif (\OProv(\gn{\OProv(\gn{\mathsf{G}})}) \lif \OProv(\gn{\lfalse})) \label{inc:inp:2in:G2-6}\\
  & \qquad\text{\ref{inc:inp:2in:G2-5}, अट P2 आणि तर्कशास्त्रावरून} \notag\\
& \OProv(\gn{\mathsf{G}}) \lif
  \OProv(\gn{\OProv(\gn{\mathsf{G}})}) \label{inc:inp:2in:G2-7}\\
   & \qquad\text{अट P3 वरून} \notag\\
& \OProv(\gn{\mathsf{G}}) \lif \OProv(\gn{\lfalse}) \label{inc:inp:2in:G2-8}\\
  & \qquad \text{\ref{inc:inp:2in:G2-6}, \ref{inc:inp:2in:G2-7} व तर्कशास्त्रावरून}\notag\\
& \OCon \lif \lnot \OProv(\gn{\mathsf{G}}) \label{inc:inp:2in:G2-9}\\
  & \qquad\text{\ref{inc:inp:2in:G2-8} च्या प्रतिवर्ती विधानावरून आणि $\OCon \ident \lnot \OProv(\gn{\lfalse})$ वरून}\notag \\
& \OCon \lif \mathsf{G} \notag\\
  & \qquad\text{\ref{inc:inp:2in:G2-1}, \ref{inc:inp:2in:G2-9} व तर्कशास्त्रावरून}\notag
\end{align}
वरील तर्कशास्त्राचा
वापर केवळ विधानीय
तर्कशास्त्रातील प्राथमिक
तथ्यांपुरता आहे.
उदाहरणार्थ,
\ref{inc:inp:2in:G2-3} मध्ये
$\Proves \lnot\varphi \liff (\varphi\lif
\lfalse)$ वापरले आहे,
आणि \ref{inc:inp:2in:G2-8} मध्ये
$\varphi \lif (\psi \lif \chi), \varphi \lif \psi
\Proves \varphi \lif \chi$ वापरले
आहे. \ref{inc:inp:2in:G2-5} आणि
\ref{inc:inp:2in:G2-6} मध्ये अट~P2
वापरताना तिची पुढील
उदाहरणे वापरली आहेत:
$\OProv(\gn{\varphi \lif \psi}) \lif
(\OProv(\gn{\varphi}) \lif \OProv(\gn{\psi}))$.
पहिल्यात $\varphi \ident
\mathsf{G}$ आणि $\psi \ident
\OProv(\gn{\mathsf{G}}) \lif \lfalse$;
दुसऱ्यात $\varphi
\ident \OProv(\gn{\mathsf{G}})$
आणि $\psi \ident \lfalse$.

दुसऱ्या अपूर्णता प्रमेयाचे
अधिक अमूर्त रूप असे:

\begin{thm}
\label{inc:inp:2in:thm:second-incompleteness-gen}
$\Th{T}$ ही~$\Th{Q}$ चा
विस्तार असलेली कोणतीही
सुसंगत, स्वयंसिद्धकीकृत
उपपत्ती असू द्या.
$\OProv[\Th{T}](y)$
हे सूत्र $\Th{T}$
साठीच्या
निष्पन्नता अटी
P1--P3 पूर्ण करत असेल,
तर $\Th{T}$ हे
$\OCon[\Th{T}]$
निष्पन्न करत नाही.
\end{thm}

\begin{prob}
$\Th{PA}$ हे
$\mathsf{G}_{\Th{PA}} \lif \OCon[\Th{PA}]$
निष्पन्न करते हे
दाखवा.
\end{prob}

\begin{digress}
या कथेतून असा बोध होतो:
गणितासाठीची कोणतीही
“स्वीकारार्ह” सुसंगत
उपपत्ती स्वतःचे
सुसंगतता-विधान
निष्पन्न करू शकत
नाही. $\Th{T}$ ही
गणिताची उपपत्ती असून
तिच्यात $\Th{Q}$ आणि
हिल्बर्टचे “सांतवादी”
युक्तिवाद (त्यांचा नेमका
अर्थ काहीही असो)
समाविष्ट आहेत असे
मानू. मागील प्रमेयातील सुसंगततेची आणि निवडलेल्या
सिद्धतायोग्यता विधेयाची सर्व अटीही पूर्ण होतात असे
गृहीत धरा. मग संपूर्ण
$\Th{T}$ हे
$\Th{T}$ चे
सुसंगतता-विधान
निष्पन्न करू
शकत नाही. त्यामुळे
तिचा सांतवादी भागही
$\Th{T}$ चे
सुसंगतता-विधान
निष्पन्न करू
शकत नाही, हे तर
अधिकच स्पष्ट आहे.
त्या अर्थाने
“सर्व गणितासाठी”
सुसंगततेची सांतवादी
सिद्धता असू शकत नाही.

“सांतवादी” या संज्ञेचा
अर्थ लावण्यात थोडी
मोकळीक आहे. गोडेलने
1931 च्या शोधपत्रात
हिल्बर्टला आकारिक
करायचे असलेल्या
गणिताबाहेरही आपण
“सांतवादी” म्हणू
असे काही असू शकेल,
ही शक्यता मान्य केली.
पण ही गोडेलची
उदार भूमिका होती:
आज सांतवादी म्हणता
येईल पण, उदाहरणार्थ,
निवडीच्या स्वयंसिद्धकासह
झर्मेलो--फ्रेंकेल संच
उपपत्ती $\Th{ZFC}$
मध्ये आकारिक करता
येणार नाही असे
काही सापडेल, हे
कल्पना करणे कठीण आहे.
\end{digress}












\section{लोबचे प्रमेय}\label{inc:inp:lob:sec}

उपपत्ती~$\Th{T}$ साठीचे गोडेल-वाक्य
हा $\lnot \OProv[\Th{T}](y)$ चा
स्थिरबिंदू असतो; म्हणजे असे
वाक्य~$\mathsf{G}$ असते की
\[\Th{T} \Proves \lnot \OProv[\Th{T}](\gn{\mathsf{G}}) \liff \mathsf{G}.
\]
सुसंगत उपपत्तीत ते निष्पन्न करता येण्याजोगे नसते. कारण
$\Th{T} \Proves \mathsf{G}$ असेल, तर
(a) निष्पन्नता अट~(1) वरून
$\Th{T} \Proves \OProv[\Th{T}](\gn{\mathsf{G}})$,
आणि (b) $\Th{T}
\Proves \mathsf{G}$ व $\Th{T} \Proves \lnot \OProv[\Th{T}](\gn{\mathsf{G}})
\liff \mathsf{G}$ यांवरून
$\Th{T} \Proves \lnot \OProv[\Th{T}](\gn{\mathsf{G}})$
मिळते. त्यामुळे $\Th{T}$ विसंगत
ठरेल. आता $\OProv[\Th{T}](y)$
च्या स्थिरबिंदूची स्थिती विचारणे
स्वाभाविक आहे. म्हणजे असे
वाक्य~$\mathsf{H}$ घेऊ की
\[\Th{T} \Proves \OProv[\Th{T}](\gn{\mathsf{H}}) \liff \mathsf{H}.
\]
ते निष्पन्न करता येण्याजोगे असेल, तर
अट~(1) वरून
$\Th{T} \Proves \OProv[\Th{T}](\gn{\mathsf{H}})$.
पण वरील सममूल्यतेला
निष्कर्षणाचा नियम लावूनही
हाच निष्कर्ष मिळतो.
म्हणून गोडेल-वाक्याच्या
बाबतीतला युक्तिवाद येथे
$\Th{T}$ विसंगत आहे
असे दाखवत नाही.
अर्थात यावरून
$\Th{T}$ हे
$\mathsf{H}$ ला \emph{खरोखर}
निष्पन्न करते असेही
सिद्ध होत नाही.

हा प्रश्न थोडा व्यापक
केल्यास पुढे जाता येते.
स्थिरबिंदूच्या
सममूल्यतेची
डावीकडून उजवीकडची
दिशा
$\OProv[\Th{T}](\gn{\mathsf{H}}) \lif \mathsf{H}$
ही \emph{प्रतिबिंबन तत्त्व}
म्हणवणाऱ्या सर्वसाधारण
रूपबंधाचे एक उदाहरण आहे:
$\OProv[\Th{T}](\gn{\varphi}) \lif \varphi$.
त्याला असे नाव दिले
आहे कारण ते एका
अर्थाने $\Th{T}$
स्वतः काय निष्पन्न
करू शकते यावर
“प्रतिबिंबन” करू शकते
असे व्यक्त करते.
मूलतः ते कोणत्याही~$\varphi$
साठी “$\Th{T}$
$\varphi$ ला निष्पन्न
करू शकत असेल, तर
$\varphi$ सत्य आहे”
असे म्हणते. हे
अर्थात फक्त सत्यनिष्ठ
उपपत्त्यांसाठीच खरे
असते. त्यामुळे
उपपत्त्या सर्वसाधारणपणे
त्याचे प्रत्येक उदाहरण
निष्पन्न करणार नाहीत
असे वाटते. मग
पुरेशी प्रबळ आणि
निष्पन्नता अटी
पूर्ण करणारी उपपत्ती
कोणती उदाहरणे
निष्पन्न करू शकते?
$\varphi$ स्वतः
निष्पन्न करता येण्याजोगे असेल
अशी सर्व उदाहरणे
नक्कीच. पुढील
निकाल दाखवतो की
एवढीच उदाहरणे शक्य
आहेत.

\begin{thm}\label{inc:inp:lob:thm:lob}
$\Th{T}$ ही~$\Th{Q}$
चा विस्तार असलेली
स्वयंसिद्धकीकरणीय
उपपत्ती असू द्या,
आणि $\OProv[\Th{T}](y)$
हे सूत्र \ref{inc:inp:prc:sec}
मधील P1--P3 अटी
पूर्ण करत असेल.
$\Th{T}$ हे
$\OProv[\Th{T}](\gn{\varphi}) \lif \varphi$
निष्पन्न करत
असेल, तर प्रत्यक्षात
$\Th{T}$ हे
$\varphi$ सुद्धा
निष्पन्न करते.
\end{thm}

दुसऱ्या शब्दांत,
$\Th{T} \Proves/ \varphi$
असेल, तर
$\Th{T} \Proves/
\OProv[\Th{T}](\gn{\varphi}) \lif \varphi$.
हा निकाल लोबचे
प्रमेय म्हणून ओळखला
जातो.

\begin{explain}
लोबच्या प्रमेयामागील
युक्ती समजण्यासाठी
सांताक्लॉज अस्तित्वात
आहे हे “सिद्ध”
करणारा चतुर युक्तिवाद
पाहा. (हा निष्कर्ष
आवडत नसेल, तर
तुम्हाला हवा असलेला
दुसरा निष्कर्ष
त्याच्या जागी
घेता येईल.)
युक्तिवाद असा:
\begin{enumerate}
\item $X$ हे “$X$
  सत्य असेल, तर
  सांताक्लॉज
  अस्तित्वात आहे”
  हे वाक्य मानू.
\item $X$ सत्य आहे
  असे धरू.
\item मग त्यात
  म्हटलेले खरे आहे:
  $X$ सत्य असेल,
  तर सांताक्लॉज
  अस्तित्वात आहे.
\item आपण $X$ सत्य
  आहे असे धरलेच
  आहे; म्हणून
  (2) आणि~(3)
  वर निष्कर्षणाचा
  नियम लावून
  सांताक्लॉज
  अस्तित्वात आहे
  असा निष्कर्ष
  काढू शकतो.
\item (2) मधील
  “$X$ सत्य आहे”
  या गृहीतकावरून
  (4) मधील
  “सांताक्लॉज
  अस्तित्वात आहे”
  हे आपण निष्पन्न
  केले. अटीची
  सिद्धता वापरून
  “$X$ सत्य असेल,
  तर सांताक्लॉज
  अस्तित्वात आहे”
  हे दाखवले.
\item पण हेच
  वाक्य~$X$ आहे.
  म्हणजे आपण
  $X$ सत्य आहे
  हे दाखवले.
\item मग वरील
  (2)--(4) च्या
  युक्तिवादाने
  सांताक्लॉज
  अस्तित्वात आहे.
\end{enumerate}
या कल्पनेचे
आकारिकीकरण करताना
“सत्य आहे” ऐवजी
“निष्पन्न करता येण्याजोगे आहे”
आणि “सांताक्लॉज
अस्तित्वात आहे”
ऐवजी~$\varphi$
घेतल्यास लोबच्या
प्रमेयाची सिद्धता
मिळते. युक्ती अशी:
स्थिरबिंदू उपपत्ती
सूत्र~
$\OProv[\Th{T}](y) \lif \varphi$
याला लावायची.
त्याचा स्थिरबिंदू
वरील रूपरेषेतील
वाक्य~$X$
याच्याशी जुळतो.
\end{explain}

\begin{proof}[याची सिद्धता: \ref{inc:inp:lob:thm:lob}]
$\varphi$ हे असे
वाक्य
असू द्या की
$\Th{T}$ हे
$\OProv[\Th{T}](\gn{\varphi}) \lif \varphi$
निष्पन्न करते.
$\psi(y)$ हे
सूत्र~
$\OProv[\Th{T}](y)
\lif \varphi$ मानू.
स्थिरबिंदू उपपत्ती
वापरून असे
वाक्य~$\theta$
शोधू की
$\Th{T}$ हे
$\theta \liff \psi(\gn{\theta})$
निष्पन्न करते.
मग पुढीलपैकी
प्रत्येक विधान
$\Th{T}$ मध्ये
निष्पन्न करता येण्याजोगे आहे:
\begin{align}
  & \theta \liff (\OProv[\Th{T}](\gn{\theta}) \lif \varphi) \label{inc:inp:lob:L-1}\\
  & \qquad \text{$\theta$ हा~$\psi(y)$ चा स्थिरबिंदू आहे}\notag \\
  & \theta \lif (\OProv[\Th{T}](\gn{\theta}) \lif \varphi) \label{inc:inp:lob:L-2}\\
  & \qquad\text{\ref{inc:inp:lob:L-1} वरून}\notag\\
  & \OProv[\Th{T}](\gn{\theta \lif (\OProv[\Th{T}](\gn{\theta}) \lif \varphi)}) \label{inc:inp:lob:L-3}\\
  & \qquad \text{\ref{inc:inp:lob:L-2} आणि अट P1 वरून}\notag \\
  & \OProv[\Th{T}](\gn{\theta}) \lif \OProv[\Th{T}](\gn{\OProv[\Th{T}](\gn{\theta}) \lif \varphi})
  \label{inc:inp:lob:L-4}\\
  &\qquad \text{\ref{inc:inp:lob:L-3} व अट P2 वरून}\notag \\
  & \OProv[\Th{T}](\gn{\theta}) \lif (\OProv[\Th{T}](\gn{\OProv[\Th{T}](\gn{\theta})}) \lif \OProv[\Th{T}](\gn{\varphi})) \label{inc:inp:lob:L-5}\\
  &\qquad \text{\ref{inc:inp:lob:L-4} व अट P2 पुन्हा वापरून} \notag\\
& \OProv[\Th{T}](\gn{\theta}) \lif \OProv[\Th{T}](\gn{\OProv[\Th{T}](\gn{\theta})}) \label{inc:inp:lob:L-6}\\
  & \qquad\text{निष्पन्नता अट P3 वरून} \notag\\
  & \OProv[\Th{T}](\gn{\theta}) \lif \OProv[\Th{T}](\gn{\varphi}) \label{inc:inp:lob:L-7} \\
  &\qquad\text{\ref{inc:inp:lob:L-5} आणि \ref{inc:inp:lob:L-6} वरून}\notag\\
  & \OProv[\Th{T}](\gn{\varphi}) \lif \varphi \label{inc:inp:lob:L-8}\\
  &\qquad\text{प्रमेयाच्या गृहीतकावरून} \notag\\
  & \OProv[\Th{T}](\gn{\theta}) \lif \varphi \label{inc:inp:lob:L-9}\\
  &\qquad\text{\ref{inc:inp:lob:L-7} आणि \ref{inc:inp:lob:L-8} वरून}\notag\\
  & (\OProv[\Th{T}](\gn{\theta}) \lif \varphi) \lif \theta \label{inc:inp:lob:L-10}\\
  & \qquad \text{\ref{inc:inp:lob:L-1} वरून}\notag \\
  & \theta \label{inc:inp:lob:L-11}\\
  & \qquad\text{\ref{inc:inp:lob:L-9} आणि \ref{inc:inp:lob:L-10} वरून}\notag \\
  & \OProv[\Th{T}](\gn{\theta}) \label{inc:inp:lob:L-12}\\
  & \qquad\text{\ref{inc:inp:lob:L-11} आणि अट~P1 वरून}\notag \\
  & \varphi \qquad\qquad\text{\ref{inc:inp:lob:L-9} आणि \ref{inc:inp:lob:L-12} वरून}\notag
\end{align}
\end{proof}

लोबचे प्रमेय मिळाल्यावर
दुसऱ्या अपूर्णता प्रमेयाची
लहान सिद्धता मिळते
(P1--P3 अटी पूर्ण
करणारे निष्पन्नता
विधेय असलेल्या उपपत्त्यांसाठी):
$\Th{T} \Proves
\OProv[\Th{T}](\gn{\lfalse}) \lif \lfalse$
असेल, तर
$\Th{T} \Proves \lfalse$.
$\Th{T}$ सुसंगत
असेल, तर
$\Th{T} \Proves/ \lfalse$.
म्हणून
$\Th{T}
\Proves/ \OProv[\Th{T}](\gn{\lfalse}) \lif \lfalse$,
म्हणजे
$\Th{T} \Proves/
\OCon[\Th{T}]$.
$\OProv[\Th{T}](x)$
चा स्थिरबिंदू~$\mathsf{H}$
हा निष्पन्न करता येण्याजोगे
आहे हेही यावरून
दाखवता येते.
कारण
\begin{align*}
  \Th{T} & \Proves \OProv[\Th{T}](\gn{\mathsf{H}}) \liff \mathsf{H}\\
  \intertext{यावरून विशेषतः}
    \Th{T} & \Proves \OProv[\Th{T}](\gn{\mathsf{H}}) \lif \mathsf{H}
\end{align*}
आणि म्हणून लोबच्या
प्रमेयावरून
$\Th{T} \Proves \mathsf{H}$.





\begin{prob}
$\Th{T}$ ही संगणनक्षमपणे
स्वयंसिद्धकीकृत उपपत्ती
असू द्या, आणि
$\OProv[\Th{T}]$ हे
$\Th{T}$ साठीचे
निष्पन्नता
विधेय असू द्या.
पुढील चार विधाने
विचारात घ्या:
\begin{enumerate}
\item जर $T \Proves \varphi$, तर $T \Proves \OProv[\Th{T}](\gn{\varphi})$.
\item $T \Proves \varphi \lif \OProv[\Th{T}](\gn{\varphi})$.
\item जर $T \Proves \OProv[\Th{T}](\gn{\varphi})$, तर $T \Proves \varphi$.
\item $T \Proves \OProv[\Th{T}](\gn{\varphi}) \lif \varphi$
\end{enumerate}
ही प्रत्येक विधाने
कोणत्या अटींवर
खरी असतात?
\end{prob}












\section{सत्याची अपरिभाष्यता}\label{inc:inp:tar:sec}

\emph{परिभाष्यता} ही संकल्पना
अंकगणिताच्या भाषेसाठी
आकारिक चिन्हार्थमीमांसा
उपलब्ध असण्यावर
अवलंबून असते.
अंकगणिताच्या भाषेतील
सूत्रे आणि वाक्ये
यांचा एक संच आपण
वर्णिला आहे.
ही वाक्ये नैसर्गिक
संख्यांबद्दल विधाने
करतात असे वाचणे
हे त्यांचे “अभिप्रेत
अर्थनिर्धारण” आहे;
असे विधान सत्य
किंवा असत्य असू
शकते.
$\Struct{N}$ ही
$\Nat$ आधारक्षेत्र
असलेली आणि
अंकगणिताच्या भाषेतील
चिन्हांचे मानक
अर्थनिर्धारण करणारी
संरचना असू द्या.
मग $\Sat{N}{\varphi}$
याचा अर्थ “$\varphi$
मानक अर्थनिर्धारणात
सत्य आहे” असा होतो.

\begin{defn}
नैसर्गिक संख्यांवरील
$R(x_1,\dots,x_k)$ हा
संबंध $\Struct{N}$
मध्ये \emph{परिभाष्य}
असतो, जर आणि फक्त जर
अंकगणिताच्या भाषेत
असे $\varphi(x_1,\dots,x_k)$
सूत्र असते की
प्रत्येक $n_1,\dots,n_k$
साठी $R(n_1,\dots,n_k)$
हे $\Sat{N}{\varphi(\num n_1,\dots,\num
  n_k)}$ असेल
तर आणि तरच खरे असते.
\end{defn}

दुसऱ्या शब्दांत,
संबंध $\Struct{N}$
मध्ये परिभाष्य
असतो, जर आणि फक्त
जर तो $\Th{TA}$
या उपपत्तीत निरूपणीय
असतो; येथे
$\Th{TA} =
\Setabs{\varphi}{\Sat{N}{\varphi}}$
हा अंकगणितातील
सत्य वाक्यांचा
संच आहे.
(हे लगेच स्पष्ट
झाले नसेल, तर
व्याख्या पुन्हा
तपासून याची
खात्री करून घ्या.)

\begin{lem}
प्रत्येक संगणनक्षम
संबंध $\Struct{N}$
मध्ये परिभाष्य
असतो.
\end{lem}

\begin{proof}
$\Th{Q}$ मध्ये
एखाद्या संबंधाचे
निरूपण करणारे
सूत्र $\Struct{N}$
मध्ये त्याच संबंधाची
परिभाषा देते, हे
सहज तपासता येते.
\end{proof}

आता उलट विधानही
खरे आहे का, असे
विचारता येते:
$\Struct{N}$ मध्ये
परिभाष्य असलेला
प्रत्येक संबंध
संगणनक्षम असतो
का? उत्तर नाही.
उदाहरणार्थ:

\begin{lem}
थांबण्याचा संबंध
$\Struct{N}$ मध्ये
परिभाष्य असतो.
\end{lem}

\begin{proof}
क्लिनीचे प्रमाण
रूप प्रमेय सांगते
की प्रत्येक आंशिक
संगणनक्षम फलन~$f$
साठी असा निर्देशांक~$e$
असतो की सर्व
$x \in \Nat$ साठी
$f(x) =
U(\umin{s}{T(e,x,s)})$;
येथे $U$ आणि $T$
आदिम पुनरावर्ती
आणि म्हणून सर्वत्र
परिभाषित आहेत.
त्यामुळे $f(x)$
परिभाषित असते
(म्हणजे संगणन
थांबते), जर आणि
फक्त जर असा~$s$
असतो की $T(e,x,s)$
खरे असते.

आता $H$ हा
थांबण्याचा संबंध
असू द्या, म्हणजे
\[H = \Setabs{\tuple{e,x}}{\lexists[s][T(e, x, s)]}.
\]
$\theta_T$ हे
$\Struct{N}$ मध्ये
$T$ ची परिभाषा
देऊ द्या. मग
\[H = \Setabs{\tuple{e,x}}{\Sat{N}{\lexists[s][\theta_T(\num e, \num x, s)]}},
\]
म्हणून $\lexists[s][\theta_T(z, x, s)]$
हे $\Struct{N}$
मध्ये~$H$ ची
परिभाषा देते.
\end{proof}

\begin{prob}
$Q(n) \defiff n \in \Setabs{\Gn{\varphi}}{\Th{Q} \Proves \varphi}$
हा संबंध अंकगणितात
परिभाष्य आहे
हे दाखवा.
\end{prob}

मग $\Th{TA}$
चे काय? ती
अंकगणितात परिभाष्य
आहे का? म्हणजे,
$\Setabs{\Gn{\varphi}}{\Sat{N}{\varphi}}$
हा संच अंकगणितात
परिभाष्य आहे
का? तार्स्कीचे
प्रमेय याचे
नकारार्थी उत्तर
देते.

\begin{thm}
\label{inc:inp:tar:thm:tarski}
अंकगणितातील सत्य
वाक्यांचा संच
अंकगणितात परिभाष्य
नाही.
\end{thm}

\begin{proof}
$\theta(x)$ त्याची
परिभाषा देते असे
समजू; म्हणजे
$\Sat{N}{\varphi}$ हे
$\Sat{N}{\theta(\gn{\varphi})}$
असेल तर आणि
तरच खरे असते.
स्थिरबिंदू उपपत्तीवरून
असे बंद सूत्र $\varphi$
असते की
$\Th{Q} \Proves \varphi \liff \lnot \theta(\gn{\varphi})$,
आणि म्हणून
$\Sat{N}{\varphi \liff \lnot \theta(\gn{\varphi})}$.
पण मग
$\Sat{N}{\varphi}$ हे
$\Sat{N}{\lnot \theta(\gn{\varphi})}$
असेल तर आणि
तरच खरे असते;
$\theta(y)$ अंकगणितातील
सत्य विधानांचा
संच परिभाषित करते
या गृहीतकाला यामुळे
विरोध होतो.
\end{proof}

तार्स्कीने हे
विश्लेषण सत्याच्या
अधिक व्यापक
तत्त्वज्ञानात्मक
संकल्पनेला लावले.
कोणतीही भाषा $L$
दिल्यास, तार्स्कीच्या
मते $L$ साठी
सत्याची पर्याप्त
संकल्पना तिच्या
प्रत्येक वाक्य
$X$ साठी पुढील
अट पूर्ण करायला
हवी:
\begin{quote}
‘$X$’ सत्य आहे
जर आणि फक्त जर
$X$.
\end{quote}
इंग्रजीसाठी
तार्स्कीने दिलेले
वारंवार उद्धृत
केले जाणारे
उदाहरण मराठीत
असे मांडता येते:
\begin{quote}
‘बर्फ पांढरा आहे’
हे सत्य आहे
जर आणि फक्त जर
बर्फ पांढरा आहे.
\end{quote}
परंतु विकर्ण
फलनाचे निरूपण
करण्याइतकी प्रबळ
असलेली कोणतीही
भाषा आणि कोणतेही
भाषिक विधेय $T(x)$
घेतल्यास, “$X$
जर आणि फक्त जर
$T(\text{‘$X$’})$
नाही” असे
वाक्य $X$
बांधता येते.
सत्यविषयक विधेयाने
एकाच वाक्याला
सत्य आणि असत्य
दोन्ही ठरवावे,
असे आपल्याला
नको आहे.
म्हणून भाषेच्या
मर्यादेबाहेर
काही ना काही
प्रकारे पाऊल
टाकल्याशिवाय
त्या भाषेतील
सर्व वाक्यांसाठी
सत्यविषयक विधेय
निर्दिष्ट करता
येत नाही,
असा तार्स्कीने
निष्कर्ष काढला.
दुसऱ्या शब्दांत,
अशा पुरेशा प्रबळ भाषेतील
सर्व वाक्यांचे सत्यविषयक विधेय
त्याच भाषेत
परिभाषित करता
येत नाही.



\part{द्वितीय-क्रम तर्कशास्त्र}
